\section{强化学习}

\subsection{基本框架}

\begin{figure}[H]
\centering
\includegraphics[width=0.95\textwidth]{slides-images/slide-014.jpg}
\caption{强化学习框架：环境给 Agent 状态 $s_t$，Agent 选择动作 $a_t$，环境返回奖励 $r_t$ 和新状态 $s_{t+1}$\protect\footnotemark}
\end{figure}
\footnotetext{Slides 第15页。}

强化学习（RL）的核心思想是让 Agent 通过与环境的大量``试错''交互来学习最优策略：执行动作 $\rightarrow$ 获得奖励 $\rightarrow$ 调整策略，使高奖励动作出现的概率增大。

\subsection{强化学习 vs 监督学习}

\begin{figure}[H]
\centering
\includegraphics[width=0.95\textwidth]{slides-images/slide-015.jpg}
\caption{强化学习（左）与监督学习（右）的对比\protect\footnotemark}
\end{figure}
\footnotetext{Slides 第16页。}

\begin{warningbox}{强化学习与监督学习的四大关键差异}
\begin{enumerate}
    \item \textbf{环境随机性}：相同动作可能导致不同结果（如推箱子时因摩擦力不均匀旋转角度不同）
    \item \textbf{信用分配（Credit Assignment）}：奖励可能延迟到很久之后才出现（如围棋直到终局才知道胜负），需要将结果归因到早期决策
    \item \textbf{不可微分性}：物理环境通常不可微，无法直接计算 $\frac{\partial r}{\partial a}$，需要大量采样做零阶梯度估计
    \item \textbf{非稳态性}：Agent 的动作会改变后续状态的分布，与监督学习中数据点互相独立截然不同
\end{enumerate}
\end{warningbox}

\subsection{Q-Learning 与 DQN}

\begin{figure}[H]
\centering
\includegraphics[width=0.95\textwidth]{slides-images/slide-018.jpg}
\caption{DQN 架构：以游戏画面为输入，经 CNN 提取特征，输出每个动作的 Q 值\protect\footnotemark}
\end{figure}
\footnotetext{Slides 第19页。}

Q 函数 $Q(s, a)$ 衡量在状态 $s$ 下执行动作 $a$ 后的\textbf{折扣期望未来累积奖励}。学习 Q 函数后，决策变得简单：选择使 Q 值最大的动作。

$$\pi^*(s) = \arg\max_a Q(s, a)$$

DeepMind 的 DQN 将 Q 函数实例化为深度卷积网络，以 Atari 游戏画面为输入，展示了令人惊叹的结果：训练 4 小时后，Agent 在 Breakout 游戏中发现了人类都不知道的``隧道策略''——将球打穿砖墙上方，实现高效消除。

\subsection{从游戏到现实：强化学习的突破}

\begin{figure}[H]
\centering
\includegraphics[width=0.95\textwidth]{slides-images/slide-022.jpg}
\caption{AlphaGo 及其后续演进：AlphaGo $\rightarrow$ AlphaGo Zero $\rightarrow$ AlphaZero $\rightarrow$ MuZero\protect\footnotemark}
\end{figure}
\footnotetext{Slides 第23页。}

\begin{knowledgebox}{Bitter Lesson：简单方法 + 规模化}
AlphaGo 到 AlphaGo Zero 的演进体现了 Rich Sutton 提出的``苦涩教训''（Bitter Lesson）：找到与规模化最兼容的简单方法，往往比精巧复杂的方法表现更好。AlphaGo Zero 去掉了所有人类知识的初始化，纯靠自我对弈，反而超越了原始 AlphaGo。
\end{knowledgebox}

\begin{figure}[H]
\centering
\includegraphics[width=0.95\textwidth]{slides-images/slide-024.jpg}
\caption{真实物理世界中的强化学习：ETH 的四足机器人在雪地、泥泞中稳健行走（左），Unitree 机器人展示极限运动能力（右）\protect\footnotemark}
\end{figure}
\footnotetext{Slides 第25页。}

在真实物理世界中，ETH 2020 年在 Science Robotics 上发表的工作令人印象深刻：在仿真中训练的四足机器人策略，直接迁移到真实的雪地、泥泞等复杂地形中，依然表现出色。关键技术是\textbf{域随机化（Domain Randomization）}——大量随机化仿真环境的物理参数（摩擦系数、地形几何等），使真实世界成为随机化分布中的一个采样点。

\begin{importantbox}{机器人运动控制接近``已解决''}
在运动控制（locomotion）领域，强化学习 + Sim-to-Real 迁移已经是成熟的解决方案。仿真不需要完美模拟所有细节（如灌木丛、积雪），只要充分随机化关键物理参数，策略就能在真实环境中鲁棒工作。
\end{importantbox}

\subsection{强化学习的局限性}

\begin{figure}[H]
\centering
\includegraphics[width=0.95\textwidth]{slides-images/slide-027.jpg}
\caption{强化学习的关键挑战：需要大量交互、安全问题、Sim-to-Real 差距在操作任务中依然显著\protect\footnotemark}
\end{figure}
\footnotetext{Slides 第28页。}

尽管在游戏和运动控制中取得了巨大成功，强化学习在\textbf{操作}（manipulation）任务中仍面临严峻挑战：

\begin{itemize}
    \item \textbf{样本效率极低}：AlphaGo Zero 需要等价于 3000 年人类知识的计算量
    \item \textbf{安全性}：训练过程中的``奇怪行为''在真实机器人上可能导致灾难性失败
    \item \textbf{可解释性差}：出错时难以诊断和修正
    \item \textbf{操作任务中 Sim-to-Real 差距更大}：仿真中抓取成功但真实中物体滑落，这种``质变''型差异难以通过域随机化消除
\end{itemize}

\subsection{本章小结}

强化学习通过试错交互学习最优策略，在游戏和运动控制领域取得了超人表现。然而，其在操作任务中因样本效率、安全性和 Sim-to-Real 差距等问题仍受限。这催生了模型学习和模仿学习等替代方案。

%% ========== Part 4: 模型学习与基于模型的规划 ==========

